% Provenance (never printed): sources/book2/BRIEFS.md, g10-03. Ricardo's model: unit labour
% requirements, autarky prices, the pattern of trade, terms of trade, gains, relative wages.
% Real numbers: figdata/grade-10/03-ricardo-model/ (Penn World Table 11.0). Brenland and Kestria's
% hours per unit: code/shared/country/ricardo.csv (invented).
\chapter{Ricardo's Model of Comparative Advantage}\label{ch:g10:ricardo-model}

A worker in the United States produced about \usd{139680} of output in 2019, measured at American
prices; a worker in Kenya, about \usd{14470}. The American worker was paid, on average, about
\usd{82660} a year in wages and other labour income, the Kenyan worker about \usd{7810}. The ratio of
outputs is close to ten, and so is the ratio of pay. The opening chart puts 139 countries on the same
picture: where output per worker is ten times higher, pay per worker is close to ten times higher too.
% figdata: wages_selected (USA 139680, 82660; KEN 14470, 7810); wages_summary (n 139)

This chart answers a fear that is as old as foreign trade. If Kenyan workers earn a tenth of American
wages, how can American firms sell anything against them? And if American workers are more productive
at everything, how can Kenyan firms sell anything at all? Both fears rest on the same mistake, and the
model of this chapter, \omterm{prop:g10:ricardo-model:wage}{Ricardo's model} of comparative advantage, shows what it is: wages adjust to
productivity, and what decides trade is not who is better at everything, but who is \emph{relatively}
better at what.

\begin{openingchart}{Labour compensation per worker against output per worker, 139 countries, 2019,
both on times-ten scales, in 2017 dollars at purchasing-power parity. Labour compensation includes wages
and an estimate of the labour income of the self-employed. Venezuela is left out: its 2019 accounts are
unreliable.}{Penn World Table 11.0; output-side real GDP,
persons engaged and share of labour compensation in GDP}
\begin{tikzpicture}
\begin{loglogaxis}[omscatter, log basis x=10, log basis y=10, log ticks with fixed point,
  xtickten={3,...,5}, ytickten={2,...,5}, width=0.84\linewidth, height=7.2cm, xmin=1000, xmax=300000,
  ymin=100, ymax=200000, xlabel={output per worker (dollars per year)},
  ylabel={labour compensation per worker (dollars)}]
\addplot[omDef, only marks, mark=*, mark size=1.3pt, fill opacity=0.7]
  table[x=opw, y=comp] {figdata/grade-10/03-ricardo-model/out/wages_productivity.csv};
\addplot[omThm, only marks, mark=square*, mark size=1.8pt, nodes near coords, point meta=explicit symbolic,
  nodes near coords style={font=\scriptsize, anchor=south east, text=omThm!80!black}]
  table[x=opw, y=comp, meta=iso3] {figdata/grade-10/03-ricardo-model/out/wages_marked.csv};
\end{loglogaxis}
\end{tikzpicture}
\end{openingchart}

\begin{recall}
From the grades 6--9 volume: a person or a country has an \emph{absolute advantage} in a good when it
makes the good with fewer hours, and a \emph{comparative advantage} when it makes the good at a lower
opportunity cost. Two parties gain from trade, even when one is better at everything, if each
specialises in its comparative advantage and they exchange at a ratio between their opportunity costs.
From \cref{ch:g10:production-possibility-frontier}: a frontier with a constant opportunity cost is a
straight line, of slope minus that cost.
\end{recall}

\section{Two countries, two goods, labour only}\label{sec:g10:ricardo-model:setup}

\omterm{prop:g10:ricardo-model:wage}{Ricardo's model} keeps the bare minimum: two countries, two goods, one factor of production --- labour
--- and a fixed number of hours of work needed for each unit of each good in each country.

\begin{definition}[Unit labour requirement]\label{def:g10:ricardo-model:unit-labour-requirement}
The \emph{unit labour requirement}\index{unit labour requirement} of a good in a country is the number of
hours of work needed to produce one unit of the good there. It is the inverse of output per hour: a
requirement of 4 hours per roll means a quarter of a roll per hour.
\end{definition}

The invented economies of this series, Brenland and Kestria, produce cloth and grain with these
requirements:

\begin{center}
\small
\begin{tabular}{lcc}
\toprule
hours per unit & cloth (one roll) & grain (one tonne) \\
\midrule
Brenland & 4 & 2 \\
Kestria & 6 & 12 \\
\bottomrule
\end{tabular}
\end{center}
% figdata: code/shared/country/ricardo.csv (invented)

Brenland needs fewer hours for both goods: it has the absolute advantage in both. To know what each
country would do without trade, take each one alone. Let Brenland and Kestria each have \num{1200} hours of
work a week. Brenland can make at most $1200/4 = 300$ rolls of cloth or $1200/2 = 600$ tonnes of grain, and
its frontier is $4C + 2G = 1200$. Kestria can make at most $1200/6 = 200$ rolls or $1200/12 = 100$ tonnes,
and its frontier is $6C + 12G = 1200$.

\begin{omfigure}
\begin{tikzpicture}[econ, x=0.0085cm, y=0.012cm]
\begin{scope}
\draw[econ axes] (0,330) node[left, font=\small] {cloth} |- (680,0) node[below left=12pt and 0pt, font=\small] {grain};
\draw[ppf] (0,300) -- (600,0);
\foreach \y in {100,200,300} \draw (0,\y) -- (-10,\y) node[left, font=\footnotesize] {\y};
\foreach \x in {200,400,600} \draw (\x,0) -- (\x,-7) node[below, font=\footnotesize] {\x};
\node[font=\small\bfseries] at (330,280) {Brenland};
\node[font=\footnotesize, anchor=west] at (320,190) {slope $-1/2$};
\end{scope}
\begin{scope}[xshift=7.6cm]
\draw[econ axes] (0,330) node[left, font=\small] {cloth} |- (680,0) node[below left=12pt and 0pt, font=\small] {grain};
\draw[ppf] (0,200) -- (100,0);
\foreach \y in {100,200,300} \draw (0,\y) -- (-10,\y) node[left, font=\footnotesize] {\y};
\foreach \x in {200,400,600} \draw (\x,0) -- (\x,-7) node[below, font=\footnotesize] {\x};
\node[font=\small\bfseries] at (330,280) {Kestria};
\node[font=\footnotesize, anchor=west] at (110,120) {slope $-2$};
\end{scope}
\end{tikzpicture}

{\small The \omterm{def:g10:production-possibility-frontier:ppf}{production possibility frontiers} of Brenland and Kestria, \num{1200} hours of work a week each,
drawn on the same scale (cloth in rolls, grain in tonnes). Brenland's lies outside Kestria's everywhere:
it has the absolute advantage in both goods. The slopes differ: a roll of cloth costs 2 tonnes of grain in
Brenland and half a tonne in Kestria.}
\end{omfigure}

Without trade, each country consumes what it produces, and the price of one good in terms of the other
is set by the hours. In Brenland a roll of cloth takes 4 hours, which could have made 2 tonnes of grain:
if cloth sold for less than 2 tonnes of grain, nobody would weave; if it sold for more, nobody would grow
grain. With both goods made, a roll of cloth trades for exactly 2 tonnes.

\begin{definition}[Autarky price]\label{def:g10:ricardo-model:autarky-price}
A country is in \emph{autarky} when it does not trade with other countries. The \emph{autarky
price}\index{autarky price} of a good is its price in autarky, measured in units of the other good. In
\omterm{prop:g10:ricardo-model:wage}{Ricardo's model}, when a country makes both goods, the autarky price of cloth in tonnes of grain is the
ratio of the \omterm{def:g10:ricardo-model:unit-labour-requirement}{unit labour requirements}, $a_C / a_G$.
\end{definition}

Brenland's \omterm{def:g10:ricardo-model:autarky-price}{autarky price} of cloth is $4/2 = 2$ tonnes of grain; Kestria's is $6/12 = 0.5$ tonne. The
\omterm{def:g10:ricardo-model:autarky-price}{autarky price} is the opportunity cost of the grades 6--9 volume, now read as a price.

\section{Who exports what}\label{sec:g10:ricardo-model:pattern}

Open the border. Let $p$ be the price of a roll of cloth in tonnes of grain on the world market, the same
for both countries (there are no transport costs in the model). A worker in either country can earn, in
an hour, either the value of the cloth made in that hour or the value of the grain.

\begin{proposition}[The pattern of trade]\label{prop:g10:ricardo-model:pattern}
Let $a_C$, $a_G$ be one country's \omterm{def:g10:ricardo-model:unit-labour-requirement}{unit labour requirements} of cloth and grain. At a world price $p$ (tonnes
of grain per roll of cloth):
\begin{enumerate}
\item if $p > a_C/a_G$, an hour spent on cloth earns more grain than an hour spent on grain, and the country
  makes only cloth;
\item if $p < a_C/a_G$, it makes only grain;
\item if $p = a_C/a_G$, both uses of an hour earn the same, and the country may make both.
\end{enumerate}
Hence, at any world price strictly between the two \omterm{def:g10:ricardo-model:autarky-price}{autarky prices}, each country makes only the good in
which it has the comparative advantage and exports it.
\end{proposition}
\begin{proof}
An hour on cloth makes $1/a_C$ rolls, worth $p/a_C$ tonnes of grain at the world price. An hour on grain
makes $1/a_G$ tonnes. The first is larger when $p/a_C > 1/a_G$, that is when $p > a_C/a_G$; workers then
move to cloth, since nothing in the model stops them, and the country makes only cloth. The other two cases
follow the same way. If $p$ lies strictly between the two \omterm{def:g10:ricardo-model:autarky-price}{autarky prices}, it is above the lower one and below
the higher one: the country with the lower \omterm{def:g10:ricardo-model:autarky-price}{autarky price} of cloth makes only cloth, and the other only
grain. Each sells abroad part of what it alone makes.
\end{proof}

For Brenland and Kestria, any world price between 0.5 and 2 tonnes of grain per roll makes Kestria a
weaver and Brenland a farmer. Kestria, worse at both, exports cloth, the good in which its disadvantage
is \emph{smaller}: 6 hours against 4 for cloth (one and a half times as many), but 12 against 2 for grain (six times as
many).

\begin{definition}[Complete specialisation]\label{def:g10:ricardo-model:complete-specialisation}
A country is in \emph{complete specialisation}\index{complete specialisation} when it produces only one of
the goods and obtains the other entirely by trade.
\end{definition}

\section{The terms of trade and the gains}\label{sec:g10:ricardo-model:gains}

\begin{definition}[Terms of trade]\label{def:g10:ricardo-model:terms-of-trade}
A country's \emph{terms of trade}\index{terms of trade} are the price of its exports in units of its
imports: how many units of the imported good one unit of the exported good buys.
\end{definition}

Kestria exports cloth: its \omterm{def:g10:ricardo-model:terms-of-trade}{terms of trade} are $p$ tonnes of grain per roll. Brenland exports grain: its
\omterm{def:g10:ricardo-model:terms-of-trade}{terms of trade} are $1/p$ rolls per tonne. A rise in $p$ improves Kestria's \omterm{def:g10:ricardo-model:terms-of-trade}{terms of trade} and worsens
Brenland's.

\begin{proposition}[Both countries gain]\label{prop:g10:ricardo-model:gains}
At any world price strictly between the two \omterm{def:g10:ricardo-model:autarky-price}{autarky prices}, each country obtains its imported good with
fewer hours of work than it would need to make that good itself.
\end{proposition}
\begin{proof}
Take Kestria, which imports grain. Making a tonne of grain itself takes $a_G^K$ hours. By trade, it makes
$1/p$ rolls of cloth, which take $a_C^K / p$ hours, and sells them for one tonne. Since $p > a_C^K/a_G^K$
(the world price is above Kestria's \omterm{def:g10:ricardo-model:autarky-price}{autarky price}), $a_C^K/p < a_G^K$: trade takes fewer hours. The
argument for Brenland, which imports cloth, is the same with the inequality reversed.
\end{proof}

\begin{example}[The gains at one roll for one tonne]\label{ex:g10:ricardo-model:gains}
Let $p = 1$. Brenland gets a roll of cloth by growing one tonne of grain, 2 hours of work, instead of
weaving it in 4: it saves 2 hours on every roll it consumes. Kestria gets a tonne of grain by weaving one
roll, 6 hours, instead of growing it in 12: it saves 6 hours on every tonne. Both gain, and the less
productive country gains more per unit at this price.
\end{example}

The gains can be drawn. A country in \omterm{def:g10:ricardo-model:complete-specialisation}{complete specialisation} can trade its whole output at the world
price, and so reach combinations it could never produce.

\begin{definition}[Consumption possibility frontier]\label{def:g10:ricardo-model:cpf}
A country's \emph{consumption possibility frontier}\index{consumption possibility frontier} is the set of
combinations of the two goods it can consume when it produces at a point of its \omterm{def:g10:production-possibility-frontier:ppf}{production possibility
frontier} and trades at the world price.
\end{definition}

Brenland, specialised in grain, makes 600 tonnes. At $p = 1$ it can exchange any part of them for cloth,
one tonne for one roll, so it can consume any $(C, G)$ with $C + G = 600$. At 150 rolls of cloth it keeps
450 tonnes of grain; without trade, 150 rolls would have taken 600 hours and left only 300 tonnes.

\begin{omfigure}
\begin{tikzpicture}[econ, x=0.011cm, y=0.011cm]
\draw[econ axes] (0,660) node[left, font=\small] {cloth (rolls)} |- (680,0) node[below left=12pt and 0pt, font=\small] {grain (tonnes)};
\draw[ppf] (0,300) -- (600,0);
\draw[curve, omThm] (0,600) -- (600,0);
\foreach \y in {150,300,600} \draw (0,\y) -- (-10,\y) node[left, font=\footnotesize] {\y};
\foreach \x in {300,450,600} \draw (\x,0) -- (\x,-10) node[below, font=\footnotesize] {\x};
\draw[guide] (0,150) -- (450,150) -- (450,0);
\draw[guide] (300,0) -- (300,150);
\node[eqpoint] at (450,150) {}; \node[above right, font=\small] at (450,150) {$T$};
\node[eqpoint, fill=omMeth] at (300,150) {}; \node[above left, font=\small] at (300,150) {$A$};
\node[eqpoint] at (600,0) {}; \node[above right, font=\small] at (600,0) {$S$};
\node[font=\footnotesize, text=omThm!85!black, anchor=west] at (180,450) {consumption possibilities, $C + G = 600$};
\draw[omMeth!60] (200,200) -- (380,265) node[font=\footnotesize, text=omMeth!85!black, anchor=west] {production possibilities};
\end{tikzpicture}

{\small Brenland at $p = 1$. It produces at $S$ (600 tonnes of grain, no cloth) and trades along the
\omterm{def:g10:ricardo-model:cpf}{consumption possibility frontier}. With 150 rolls of cloth it consumes 450 tonnes of grain at $T$; on its
own it would have only 300 at $A$.}
\end{omfigure}

Where, between 0.5 and 2, does the world price settle? The model so far says only that it lies in that
range. Inside the range, each country supplies the world with all of its one good, so the world makes
200 rolls of cloth (Kestria) and 600 tonnes of grain (Brenland): a relative supply of $200/600 = 1/3$ roll per
tonne. At $p = 0.5$ Kestria is willing to make any mix; at $p = 2$ Brenland is. The world relative supply of
cloth is therefore a staircase. The price is set where this staircase meets the world's relative demand for
cloth, which depends on what people want to buy, the subject of the next chapter.

\begin{omfigure}
\begin{tikzpicture}[econ, x=3.2cm, y=1.6cm]
\draw[econ axes] (0,2.7) node[left, font=\small, align=right] {$p$ (tonnes\\per roll)} |- (2.0,0) node[below left=12pt and 0pt, font=\small] {world cloth $\div$ world grain};
\draw[supply] (0,0.5) -- (1/3,0.5) -- (1/3,2) -- (1.85,2);
\draw[demand] (0.08,2.55) .. controls (0.25,1.3) and (0.5,0.7) .. (1.6,0.25);
\foreach \y/\l in {0.5/0.5, 1/1, 2/2} \draw (0,\y) -- (-0.04,\y) node[left, font=\footnotesize] {\l};
\draw (1/3,0) -- (1/3,-0.06) node[below, font=\footnotesize] {$1/3$};
\node[font=\footnotesize, text=omThm!85!black, anchor=west] at (1.0,2.15) {relative supply};
\node[font=\footnotesize, text=omDef!85!black, anchor=west] at (0.62,1.3) {relative demand};
\draw[guide] (0,1) -- (1/3,1);
\node[eqpoint] at (1/3,1) {};
\end{tikzpicture}

{\small The world relative supply of cloth is a staircase: no cloth below $p = 0.5$, Kestria's 200 rolls
against Brenland's 600 tonnes between 0.5 and 2, and any amount of cloth from Brenland too at $p = 2$. A
relative demand curve (drawn here so that it crosses at $p = 1$) sets the price.}
\end{omfigure}

\section{Wages follow productivity}\label{sec:g10:ricardo-model:wages}

Trade is done by firms, which compare money costs, not hours. Money costs depend on wages. \omterm{prop:g10:ricardo-model:wage}{Ricardo's model}
says where wages end up.

In \omterm{def:g10:ricardo-model:complete-specialisation}{complete specialisation}, a Brenland worker makes half a tonne of grain an hour, and a Kestria worker a
sixth of a roll of cloth. With competing employers and free movement of workers inside each country, the
hourly wage equals the value of what a worker makes in an hour. Measure prices in tonnes of grain: a tonne
of grain is worth 1 and a roll of cloth $p$. Then
\[
  w_B = \tfrac12, \qquad w_K = \tfrac{p}{6}, \qquad \frac{w_B}{w_K} = \frac{3}{p}.
\]

\begin{proposition}[The relative wage]\label{prop:g10:ricardo-model:wage}
In \omterm{def:g10:ricardo-model:complete-specialisation}{complete specialisation}, the ratio of Brenland's wage to Kestria's lies between the two ratios of output
per hour: Brenland's output per hour is $6/4 = 1.5$ times Kestria's in cloth and $12/2 = 6$ times in grain,
and $1.5 < w_B/w_K < 6$.
\end{proposition}
\begin{proof}
$w_B/w_K = 3/p$, and $0.5 < p < 2$, so $3/2 < 3/p < 3/0.5$, that is $1.5 < w_B/w_K < 6$.
\end{proof}

At $p = 1$, Brenland's wage is three times Kestria's. Compare money costs per unit, in units of Kestria's
hourly wage: a roll of cloth costs 6 in Kestria and $4 \times 3 = 12$ in Brenland, so Kestria's firms sell
cloth more cheaply; a tonne of grain costs 12 in Kestria and $2 \times 3 = 6$ in Brenland, so Brenland's firms
sell grain more cheaply. Each country undersells the other in one good. Kestria's low wage does not let it
undersell Brenland in everything, because it is low exactly as far as Kestria's productivity is low; and
Brenland's high wage does not price it out of everything, because it is high only up to its productivity
lead in grain.

The opening chart is consistent with this prediction across countries. Across 139 countries, the line that
best fits the cloud on the two times-ten scales has a slope of 1.01: ten times the output per worker goes
with about ten times the pay. The cloud is tight: the correlation of the two (on the times-ten scales) is
0.96. It is not a test, because pay is measured as a share of output: with similar labour shares, the link
is close to built in.
% figdata: wages_summary (n 139, slope 1.01, r 0.96)

\begin{measure}[Output per worker across countries]
Output per worker is GDP divided by the number of persons engaged (employees and the self-employed). To
compare countries, GDP is converted not at the market exchange rate but at \emph{purchasing-power
parities}: rates that make a common basket of goods cost the same everywhere, so that a haircut in Nairobi
counts as much as a haircut in Chicago. Labour compensation is GDP times the share of it paid to labour; for
the self-employed, whose income mixes pay for work and return on their capital, the database estimates the
labour part. Both axes of the opening chart come from the same national accounts, so part of their link is
built in: what the chart adds is that the labour share is roughly similar across countries, so pay is
roughly proportional to output.
\end{measure}

\begin{model}{Ricardo's model}\index{Ricardo's model}
\assumes Two countries, two goods; labour is the only factor of production; \omterm{def:g10:ricardo-model:unit-labour-requirement}{unit labour requirements} are
  fixed; workers move freely between industries inside a country but not between countries; no transport
  costs; competition.
\predicts Each country exports the good with the lower relative labour requirement; the \omterm{def:g10:ricardo-model:terms-of-trade}{terms of trade} lie
  between the \omterm{def:g10:ricardo-model:autarky-price}{autarky prices}; both countries gain; the relative wage lies between the two ratios of output
  per hour.
\tested Japan's opening to trade in 1859, told in the grades 6--9 volume, gives a direct test: valued at
  the prices Japan had before it opened, its net exports had the negative sign that comparative advantage
  predicts in every year from 1868 to 1875 (Bernhofen and Brown 2004). That pay follows output per worker
  across 139 countries (opening chart) is consistent with the wage prediction, not a test of it.
% ledger: g7:bernhofen-brown-2004
\breaks Other factors (land, capital, skills) matter and are owned by different people, so some people inside a
  country lose from trade even when the country gains (grade 12); \omterm{def:g10:ricardo-model:complete-specialisation}{complete specialisation} is rare, because
  costs rise as an industry grows; transport costs and tariffs keep some goods out of trade altogether.
\end{model}

The model's lesson about the losers deserves a sentence now. Brenland's weavers, when the border opens, must
move to the farms. In the model this costs nothing; in reality it can cost years of income. The country's
gain is real, and so is the weavers' loss; grade 12 measures both.

\section{Exercises}

\begin{exercise}[$\star$]\label{exo:g10:ricardo-model:1}
Using Brenland and Kestria's table of hours, say which country has the absolute advantage in cloth and which
in grain. Explain why absolute advantage alone cannot tell which country exports which good.
\end{exercise}

\begin{exercise}[$\star$]\label{exo:g10:ricardo-model:2}
Compute the \omterm{def:g10:ricardo-model:autarky-price}{autarky price} of grain, in rolls of cloth per tonne, in Brenland and in Kestria. Check that each is
the inverse of the \omterm{def:g10:ricardo-model:autarky-price}{autarky price} of cloth given in the chapter.
\end{exercise}

\begin{exercise}[$\star$]\label{exo:g10:ricardo-model:3}
Country A makes a bicycle in 20 hours and a chair in 5. Country B makes a bicycle in 30 hours and a chair in
15. Which country has the comparative advantage in bicycles? Between which two prices (chairs per bicycle) must
the world price lie for both to trade?
\end{exercise}

\begin{exercise}[$\star$]\label{exo:g10:ricardo-model:4}
Read the opening chart. About how much labour compensation per worker corresponds to output per worker of
\usd{10000}? Of \usd{100000}?
\end{exercise}

\begin{exercise}[$\star$]\label{exo:g10:ricardo-model:5}
``Kestria cannot sell anything to Brenland, since Brenland makes everything with fewer hours.'' Explain in three
sentences why this is wrong.
\end{exercise}

\begin{exercise}[$\star\star$]\label{exo:g10:ricardo-model:6}
Let the world price be $p = 1.5$ tonnes of grain per roll. How many hours does Brenland save on each roll of
cloth it imports? How many does Kestria save on each tonne of grain? Compare with
\cref{ex:g10:ricardo-model:gains}: which country's gain rose, and why?
\end{exercise}

\begin{exercise}[$\star\star$]\label{exo:g10:ricardo-model:7}
A new plough cuts Kestria's hours per tonne of grain from 12 to 3. Predict what happens to the pattern of trade
and to the gains. What if it cut them to 2?
\end{exercise}

\begin{exercise}[$\star\star$]\label{exo:g10:ricardo-model:8}
At $p = 1.5$, what is the ratio of Brenland's wage to Kestria's? Compare the money cost of a roll of cloth and
of a tonne of grain in the two countries, in units of Kestria's hourly wage, and say which country sells each
good more cheaply.
\end{exercise}

\begin{exercise}[$\star\star$]\label{exo:g10:ricardo-model:9}
Kestria (\num{1200} hours) specialises in cloth and trades at $p = 1.5$. Write its \omterm{def:g10:ricardo-model:cpf}{consumption possibility
frontier}. How much grain can it consume if it keeps 100 rolls of cloth? How much without trade?
\end{exercise}

\begin{exercise}[$\star\star$]\label{exo:g10:ricardo-model:10}
% figdata: wages_selected (USA 139680, 82660; DEU; MEX; CHN; IND 19440, 10370; KEN)
The table behind the opening chart gives, for the United States, output per worker \usd{139680} and
compensation \usd{82660}; for India, \usd{19440} and \usd{10370}. Compute the ratio of output per worker and the
ratio of compensation, United States to India, and the share of output paid to labour in each country.
\end{exercise}

\begin{exercise}[$\star\star$]\label{exo:g10:ricardo-model:11}
Without trade, Brenland and Kestria each split their \num{1200} hours half and half between the two goods. Compute
the world output of each good. Then let Kestria make only cloth while Brenland makes the 50 rolls of cloth still
needed to keep world cloth at its \omterm{def:g10:ricardo-model:autarky-price}{autarky} level, and grain with the rest of its hours. By how much does world
grain rise?
\end{exercise}

\begin{exercise}[$\star\star\star$]\label{exo:g10:ricardo-model:12}
Describe an observation about trade between two countries, with their \omterm{def:g10:ricardo-model:unit-labour-requirement}{unit labour requirements} known, that would
contradict \omterm{prop:g10:ricardo-model:wage}{Ricardo's model}. Then give a reason, outside the model, why such an observation might still be found.
\end{exercise}

\begin{exercise}[$\star\star\star$]\label{exo:g10:ricardo-model:13}
Brenland and Kestria make three goods. Hours per unit (Brenland, Kestria): cloth 4 and 6, grain 2 and 12, tools 5
and 5. Rank the goods by Kestria's hours divided by Brenland's. A good is made where its money cost, hours times
the wage, is lower. Show that if Kestria made grain more cheaply than Brenland, it would make all three goods more
cheaply, and explain why trade, which must be paid for, rules this out. Show that if Brenland makes cloth more
cheaply, it also makes grain more cheaply. If the wage ratio $w_B/w_K$ lies between 1 and 1.5, which country makes
each good?
\end{exercise}

\begin{exercise}[$\star\star\star$]\label{exo:g10:ricardo-model:14}
Explain, with the numbers of the chapter at $p = 1$, why Brenland, whose wage is three times Kestria's, still
exports grain. What would happen to Brenland's exports if its wage were seven times Kestria's?
\end{exercise}

\begin{exercise}[$\star\star\star$]\label{exo:g10:ricardo-model:15}
Prove that two countries trade, each exporting one good, only if the world price lies between their two \omterm{def:g10:ricardo-model:autarky-price}{autarky
prices} (or equals one of them). Use \cref{prop:g10:ricardo-model:pattern}: what does each country make if the
price is above both \omterm{def:g10:ricardo-model:autarky-price}{autarky prices}?
\end{exercise}

\section{Problem: Opening the Bren--Kest Border}

\begin{problem}[{Weekend problem --- autarky, specialisation, trade at one for one, and who gains how much}]
\label{pb:g10:ricardo-model:1}
Brenland has 600 million hours of work a year and Kestria \num{1800} million (it has more workers, each less
productive). The hours per unit are those of the chapter: cloth 4 and grain 2 in Brenland, cloth 6 and grain 12 in
Kestria. In \omterm{def:g10:ricardo-model:autarky-price}{autarky}, each country spends half of its hours on each good. Quantities are in millions of rolls and
millions of tonnes.

\medskip
\noindent\textbf{Part I --- Autarky.}
\begin{enumerate}
\item Write each country's \omterm{def:g10:production-possibility-frontier:ppf}{production possibility frontier} and its two ends.
\item Give each country's \omterm{def:g10:ricardo-model:autarky-price}{autarky price} of cloth, in tonnes of grain per roll.
\item Compute what each country produces and consumes in \omterm{def:g10:ricardo-model:autarky-price}{autarky}.
\item Compute world output of cloth and of grain in \omterm{def:g10:ricardo-model:autarky-price}{autarky}.
\end{enumerate}

\medskip
\noindent\textbf{Part II --- The pattern.}
\begin{enumerate}[resume]
\item Which country has the absolute advantage in each good?
\item Which country has the comparative advantage in each good?
\item Between which two values must the world price of cloth lie for both countries to specialise?
\end{enumerate}

\medskip
\noindent\textbf{Part III --- Trade at one roll for one tonne.}
\begin{enumerate}[resume]
\item The border opens and the world price settles at $p = 1$. What does each country produce?
\item Each country still spends half the value of its output on each good. What does each consume?
\item What does each country export and import? Check that world markets clear.
\item Compare world output with \omterm{def:g10:ricardo-model:autarky-price}{autarky}.
\item Compute each country's gain in goods, compared with \omterm{def:g10:ricardo-model:autarky-price}{autarky}.
\item How many hours does Brenland save on each roll of cloth it consumes? How many does Kestria save on each tonne of
  grain?
\end{enumerate}

\medskip
\noindent\textbf{Part IV --- Wages and a new plough.}
\begin{enumerate}[resume]
\item Measuring prices in tonnes of grain, compute each country's hourly wage and their ratio.
\item Check that, at these wages, Kestria makes cloth more cheaply than Brenland, and Brenland grain more cheaply than
  Kestria.
\item Between which two values can the relative wage lie, whatever the world price in the range of question 7?
\item A new plough cuts Kestria's hours per tonne of grain from 12 to 8. Does the pattern of trade change? Is $p = 1$
  still in the range where both specialise?
\item With the plough and $p = 1$, how many hours does Kestria now save per tonne of grain imported? Explain why a better
  plough reduces Kestria's gain from importing grain.
\item At $p = 1$, what plough (hours per tonne of grain in Kestria) would leave Kestria with no gain from
  importing grain?
\item State the result of the problem: each country's gain from trade at $p = 1$ (before the plough), in hours of work
  saved per unit of its imported good, and the relative wage.
\end{enumerate}
\end{problem}
